\documentclass[11pt]{article}
\usepackage[margin=1in]{geometry}
\usepackage{tgtermes}
\usepackage{setspace}
\setstretch{1.15}
\usepackage{microtype}
\usepackage{amsmath,amssymb}
\usepackage[hidelinks,pdftitle={Two modalities in one quantifier},pdfauthor={}]{hyperref}
\usepackage[authoryear,round]{natbib}
\setcitestyle{aysep={}}
\pagestyle{plain}
\begin{document}
\begin{center}
{\Large Two modalities in one quantifier\par}
\vspace{0.7em}
\end{center}

\begin{abstract}
A process function is defined by a universal quantifier over local maps and an existence-and-uniqueness claim about histories. Those are different modal claims. A local map can lie in the range of the quantifier and still leave the fixed-point equation with no solution, or with two. A cyclic wiring can lie in the same range and have exactly one solution. The canonical unitary biconditional says that unitarity of one closed loop, under every local unitary and every matched finite ancilla, coincides with the history claim. That coincidence is a necessary connection inside the model. It is not a de re necessity about devices, not a free recombination of events, and not a result about grounding or monism. Counting solutions settles those numbers. It does not make any tuple actual, and it does not answer an open question about infinite maps. The argument is an alternative adaptation of a public formalism. The formalism is doing the philosophical work. No further ontology is added.
\end{abstract}

\noindent\textit{Keywords:} modality; compossibility; quantifiers; process functions; necessity

\section{The distinction the formalism already draws}
\label{sec:dist}

Let a wiring be a map from output tuples to input tuples, on finite alphabets, one component per party. A local map at a party is a function from that party's input alphabet to its output alphabet. A history for a tuple of local maps is an output tuple fixed by the composition of those maps with the wiring. The wiring is a process function when every tuple of local maps has exactly one history \citep{baumeler2016}. The public theorem used below adds one biconditional. A canonical permutation of a source-sink space, defined for every wiring, yields a closed loop that is unitary for every local unitary and every matched finite ancilla if and only if the wiring is a process function. The permutation itself exists whether or not any history does. These statements are not re-proved here. They are the material.

The metaphysical reading this paper defends is small, and it is read off the quantifier rather than laid on top of it. The quantifier ``for every tuple of local maps'' treats each tuple as a case. Call that substitution-possibility: the tuple is among the cases the definition asks about. The clause ``there is exactly one history'' treats output tuples as jointly satisfying, or failing to satisfy, the fixed-point equation for that case. Call the joint satisfaction history-compossibility. Substitution-possibility and history-compossibility are not the same modality. A tuple can be a case and have no solution. A tuple can be a case and have two solutions. A tuple can be a case and have one solution even though each output depends on the others. The formalism records all three. A metaphysical vocabulary that collapses them will mis-state the formalism.

Nothing in that paragraph says that any local map is metaphysically possible in the world, that any history is a real event, or that the world is or is not a unity. Those claims are not supported by the distinction, and this paper does not make them. The substance is the distinction. If it seems too thin to be metaphysics, the reply is in Section~\ref{sec:obj}: a distinction between two modal roles inside one definition is already a modal claim, and inventing a further ontology would not make it more honest.

Other adaptations of the same mathematics ask different questions: what an intervention test would identify, whether a simulator implements a wiring, whether a robustness notion has been transferred too far. Those questions are not this one. A reader who thinks they are the same question is right that the mathematics is the same and wrong that the modal roles of the two clauses were already the topic. This paper is an alternative adaptation. It does not claim an independent discovery.

\section{Three cases, three modal profiles}
\label{sec:cases}

The cases are calculations. They are not stories about agents, and they are not evidence that a device exists.

First, two parties share a bit. One local map copies the bit. The other flips it. The tuple of maps is a case: the quantifier ranges over it. The fixed-point equation has no solution. Substitution-possibility holds and history-compossibility fails. It is a mistake to say the flip was impossible. The flip was specified. What is impossible, relative to that specification, is a pair of outputs that satisfies both maps at once. The impossibility is joint. It is not an impossibility of either map taken as a substitution.

Second, both maps copy the bit. The tuple is again a case. Both bit values are solutions. History-compossibility does not fail by emptiness. It fails by surplus, if the claim at issue is uniqueness. There are two ways for the outputs to fit the maps. A modality that says only ``a history is possible'' is too coarse. The formalism distinguishes one fitting tuple from two. Uniqueness is part of the history claim in the definition of a process function, not a decorative extra.

Third, three bits are wired by \(w(a,b,c)=(\lnot b\land c,\ \lnot c\land a,\ \lnot a\land b)\). Every tuple of binary maps, including the three copying maps, has exactly one history. For the copying maps the history is \((0,0,0)\). Each bit depends on the others. The dependence is cyclic in the ordinary sense that there is no ordering of the bits in which each depends only on earlier ones. Substitution-possibility and unique history-compossibility both hold. Cyclicity is not, in this formalism, a bar to the history claim. A metaphysical picture on which a cycle is automatically an impossible history is not this formalism's picture. The formalism allows the cycle and counts the solutions.

The three profiles are enough to separate the modalities. If substitution-possibility implied unique history-compossibility, the first two wirings would not be cases. They are cases. If unique history-compossibility implied an acyclic order, the third wiring would not be a process function. It is one. The separations are finite and checkable. They do not depend on a theory of possible worlds. They depend on counting solutions of stated equations.

\section{What the biconditional adds, and what it does not}
\label{sec:iff}

The unitary biconditional connects a third claim to the history claim. Call it loop-possibility: for every local unitary on a matched finite ancilla, the closed loop is unitary. The theorem says loop-possibility holds if and only if every substitution has exactly one history. The connection is necessary and sufficient inside the model, for that canonical construction and that class of unitaries. It is not the claim that some reversible encoding or other exists. The encoding \(U\lvert o,e\rangle=\lvert w(o)+e,o\rangle\) exists for the first wiring and the second, which are not process functions. Bare reversibility of that map is not loop-possibility. It is a permutation of a basis. Treating it as the metaphysical possibility of the wiring would erase the biconditional, because the permutation does not track the history claim.

The necessity in the biconditional is the necessity of a proved equivalence. If the history claim holds, the loop claim holds, and conversely, under the stated definitions. That is de dicto necessity of a biconditional sentence about maps and operators. It is not a de re necessity about a particular device, because no device has been identified. It is not a claim that the world could not have failed to satisfy the equivalence. The equivalence is about the objects the definitions introduce. A world in which those objects are not instantiated is not a counterexample to the equivalence. A device that fails to match the canonical loop is not a counterexample either. It is a device that is not the loop.

Two limits keep the necessity from spreading. The proof uses finite alphabets and finite-dimensional ancillas. It does not say that every exclusive self-map of an infinite product of bits has a fixed point. Whether that is true for Baire-class-one maps is open, and the openness is a reminder that existence and uniqueness already come apart when the quantifier is no longer a finite check. This paper does not close the question and does not rest a metaphysical thesis on it. The converse half of the biconditional uses one family of local unitaries, the reversible encodings of functions. It does not say that an arbitrary purification is unitary only for process functions. A modal reading that says ``reversibility is possible only if there is a unique history'' has dropped the canonical restriction and the class of tests. The formalism does not drop them.

A known distribution, computed in one basis from a known board, makes a linear functional negative, while the functional is nonnegative on three-bit process functions in an enumeration this paper does not repeat. If the enumeration is granted, that distribution is not a mixture of process functions. The modal point is narrow. Being inducible in one basis by a pure process, in the model's sense of purity, does not imply being a convex combination of the wirings that satisfy the history claim. Purity in that sense is not history-compossibility. The calculation is not a discovery about the space of metaphysical possibilities. It is a separation of two sets of mathematical objects, in a named frame. Another frame need not repeat it. A fixed basis is not a universal ontology.

\section{Necessity of the equivalence is not necessity of a history}
\label{sec:nec}

The biconditional is necessary in a philosopher's sense that is easy to overread. Given the definitions, there is no model of those definitions in which the history claim holds and the loop claim fails, or the converse, for the canonical construction. The necessity is conditional on the definitions. It does not attach to a history the way necessity attaches to a theorem's being true. A particular history, such as \((0,0,0)\) for the copying case of the three-bit rule, is the unique solution of an equation. One can say, loosely, that the history is necessary relative to those maps, meaning no other tuple solves the equation. That relative necessity is uniqueness. It is not the necessity of the biconditional. Uniqueness can hold for one substitution while the biconditional's universal claim is about all substitutions. The copying case of the third wiring is unique. The biconditional's history side requires uniqueness for every substitution, which the third wiring also has, and which a wiring with one unique case and one empty case does not.

Relative necessity of a history also fails to be absolute necessity. ``No other tuple solves the equation'' is a fact about that equation. It does not say the history obtains no matter what the local maps are. Change the maps and the solution set can change. The first wiring, with a flip in place of a copy, has an empty solution set. The history that was relatively necessary under copying is not necessary under the flip, and is not a solution at all. A metaphysical necessity that survived every substitution would be a tuple that solved every system. No output tuple does that, in these examples. The formalism therefore contains relative necessities that are fragile under substitution. Fragility under substitution is another way the two modalities come apart. The case is still a case. The solution that was unique for a different case is gone.

None of this is necessity in the sense of a law of nature, a logical truth independent of the definitions, or a metaphysical law about events. It is necessity relative to a specified equation, and necessity of an equivalence proved from definitions. Keeping the grades separate is the point of the section. A sentence that says ``the history is necessary, and the unitary law is necessary, so the device must exist'' has run three grades together and added an existence claim about a device that none of the grades contained.

The open infinite question is a fourth grade, and it should not be folded in. For finite alphabets, uniqueness and existence are both settled by counting, and the biconditional uses both. For an exclusive self-map of an infinite product, the exclusivity argument still yields at most one fixed point, while existence is not given by the same counting. The open question is whether existence follows for a stated class of maps. That question is not a hidden necessity in the finite theorem. A philosopher who cited the finite biconditional as proof that an infinite exclusive system must have a state would be citing a theorem whose hypotheses are false of that system. The finite necessity does not export.

\section{What a recombination principle would have to say}
\label{sec:rec}

It is natural to hear the universal quantifier as a recombination principle: the local maps may be chosen independently, and the wiring must still make sense. The formalism agrees with the first half and refuses a strong reading of the second. Every tuple is a case. Not every case has one history. A principle that says ``any local maps can be set together'' is a principle about substitution-possibility. A principle that says ``any local maps can be set together and there will be a unique history'' is a principle about both modalities at once. The first two wirings satisfy a principle of the first kind, in the thin sense that they are in the quantifier's range, and they refute a principle of the second kind. The refutation is not a refutation of recombination as such. It is a refutation of the slide from one formulation to the other.

The third wiring shows a different slide. One might think that if the maps recombine freely and a unique history exists, the history must be readable as an ordered sequence of independent choices. The cyclic example has a unique history and no such order. Recombination of the local maps, in the sense of the quantifier, does not produce an ordered decomposition of the history. The history is one tuple. The order is absent. A metaphysics that requires every compossible tuple to be the trace of an ordered process is adding a requirement the example does not meet.

This is not an interpretation of any particular recombination theory. No text of Lewis, or of anyone else, is being glossed. The point is a constraint on what such a theory can take this formalism to show. It can take the formalism to show that independent specifiability of maps is not unique joint satisfiability of outputs. It cannot take the formalism to show that a favored theory of possible worlds is true or false, because the formalism does not speak in that theory's terms. Attributing the constraint to a named metaphysician without a passage would be a different paper, and it is not this one.

The same restraint applies to grounding and to monism. A monist can accept the distinction between substitution-possibility and history-compossibility. A pluralist can accept it. A grounding theorist can accept it and still ask which facts ground which. The distinction does not answer that question. It does not say that histories are nothing over and above maps, or that maps are nothing over and above histories. In the first wiring there is a tuple of maps and no history. In the third there is a history that is not an ordered product of independent choices. Neither fact is a priority claim. Using them as a refutation of monism, or of grounding, would be an argument this paper does not have and does not assert.

\section{Modal scope of the finite check}
\label{sec:scope}

Because the alphabets are finite, each case can be checked by counting. The check is a decision about a mathematical question. It is not a survey of metaphysical possibilities. Two confusions follow if the check is asked to do more than it does.

The first confusion treats decidability as a modal contraction: since we can check every case, there is no room for an unrealized possibility, and the model's modal structure is exhausted by what the check prints. The check prints, for each substitution, a number of histories. That number is the history claim for that case. Substitution-possibility is not printed as a separate line, because every substitution was fed to the check. The check's ability to take every substitution as input is the substitution-possibility. The printout's number is the history claim. A procedure that can ask every question is not a procedure for which every question has one answer. The first wiring is the illustration. The check runs. The number is zero.

The second confusion treats the check as an implementation of the possibilities it counts. A machine that prints ``one'' has settled a mathematical question about a table. It has not made a history actual. Actuality is not one of the modalities in the definition. The definition does not say that the unique history obtains. It says that there is exactly one tuple that would satisfy the equation if the local maps were those. The ``if'' is inside the model. A metaphysics that hears ``there is a unique history'' as ``a particular event is actual'' has changed the modal force of the existential quantifier. The existential quantifier in the definition is a quantifier over tuples of symbols. It is not a tensed claim that an event occurred.

The unitary side has the same structure. ``The loop is unitary for every local unitary'' is a universal claim about operators. A check that multiplies matrices and prints ``unitary'' has settled that claim for the operators it was given. It has not actualized a physical process. The biconditional says the printout on the unitary side matches the printout on the history side, for the canonical construction. Matching printouts are not two actualities. They are two mathematical claims with the same truth value.

Finite search also blocks a different inflation. The open question about Baire-class-one maps is a place where one cannot simply count. The finite case is the place where one can. The contrast does not show that finite histories are metaphysically innocent and infinite ones are deep. It shows that the finite modal claims in this paper are claims one can settle by enumeration, and that settling them does not settle the infinite question. A metaphysics that used the finite biconditional as a stand-in for the infinite existence question would be using a settled case as a disguise for an open one.

\section{Objections}
\label{sec:obj}

The first objection says that a distinction between two clauses of a definition is grammar, not metaphysics. The clauses have different truth conditions, and the truth conditions come apart on examples. A pair of claims with different truth conditions is a pair of claims. Calling both ``possibility'' without saying which is the metaphysical error the paper is about. If the objector drops the modal vocabulary and says only that the definition has two clauses, the paper's conclusion can be restated: those clauses are independent in the three ways Section~\ref{sec:cases} lists, and a biconditional ties a third clause to one of them and not to the bare permutation. The restatement is still a claim about what follows from what. That is as metaphysical as the material gets. Adding substances, worlds, or grounding relations would be an invention, not a deepening.

The second objection says that substitution-possibility is not a modality at all, because the quantifier is an ordinary universal quantifier over a finite set. Ordinary universal quantifiers have modal analogues when they are used to say what is a case rather than what is a solution. The paper's term marks that use. It does not claim that the set of functions is a set of possible worlds. A reader who dislikes the term can substitute ``being in the domain of the quantifier'' throughout. The separations remain. The term is dispensable. The separations are not.

The third objection says that uniqueness is an epistemic demand, a demand to know which history, rather than a modal fact. In the second wiring the equation has two solutions whether or not anyone is ignorant. The surplus is in the solution set. Ignorance might matter for a different question, about what an agent can predict. This paper does not discuss agents. The cardinality of the solution set is a property of the equation.

The fourth objection says that the cyclic wiring shows time travel to be possible, or shows a block universe, or shows the opposite. It shows that a particular system of Boolean equations has one solution and a cyclic dependency graph. Possibility of time travel would require an identification of the equations with a physical or metaphysical process in time. No identification is given. The Mind-journal adaptation of this mathematics discusses operations and histories in a different vocabulary and is not repeated here. The present point is only that cyclicity does not, inside this formalism, refute uniqueness, and uniqueness does not supply an order. Those negatives do not add up to a cosmology.

The fifth objection says the paper is a physics note with a modal paragraph. The note's theorems are used as the examples, because the examples are exact. The argument does not consist in re-proving them. It consists in saying which modal roles their clauses occupy and which further modal claims they do not license. A reader who wants the proofs should read the technical note. A reader who wants a refutation of monism will not find it, and the absence is not an oversight. The premises do not support that conclusion. Stating a conclusion the premises do not support would be a different failure from the one this objection names.

\section{Conclusion}
\label{sec:end}

The process-function quantifier ranges over tuples of local maps. That is one modal role: each tuple is a case. The fixed-point clause says how many output tuples fit a case. That is another role. The roles come apart. Some cases have no fitting tuple. Some have two. Some have one, and among those some are cyclic. The canonical biconditional ties unitarity of a stated loop, under a stated class of unitaries, to the second role, not to the mere existence of a reversible relabeling. The tie is a necessary equivalence inside the model. It is not a de re necessity, not a recombination principle that guarantees histories, and not a thesis about grounding or monism.

The finite check settles the numbers. It does not actualize the tuples it counts, and it does not answer the open question about infinite exclusive maps. The metaphysical content of the formalism, on this adaptation, is the set of separations just listed. They are enough to stop a family of slides. They are not enough to start a theory of possible worlds. The paper stops where the separations stop.

\section*{Use of a language model}
Grok Bot, a Cursor cloud agent, drafted this manuscript on 8 November 2026. The use was generative, not confined to copy editing. There is no methods section; this note is the alternative location the journal's guidelines allow. The model version is not recorded. The agent is not an author. Verification is pending. No generative images were made.

\begin{thebibliography}{99}
\bibitem[Baumeler and Wolf(2016)]{baumeler2016} Baumeler, \"A., and S. Wolf. 2016. ``The Space of Logically Consistent Classical Processes without Causal Order.'' \textit{New Journal of Physics} 18:013036. \url{https://doi.org/10.1088/1367-2630/18/1/013036}
\end{thebibliography}
\end{document}
